\documentclass[11pt,a4paper]{article}
% preamble.tex - shared house preamble for /explain-projects documents.
% Every package here is VERIFIED present in this machine's TeX Live install.
% Do not add packages without running: kpsewhich <pkg>.sty
% Notably ABSENT (do not \usepackage them, the build will fail):
%   algorithm2e, algorithm, algorithmic, algpseudocode  (texlive-science is not installed)
%   -> use the \begin{pseudocode} tcolorbox environment defined below instead.
%   lmodern, charter, mathptmx, times, mathpazo, courier, newtx*, tgtermes
%   -> their .sty files may resolve but the FONT FILES are absent, so the build
%      dies with "I can't find file `bchr8t'" or similar. Verified 2026-07-16 by
%      test-compiling each one. Only Computer Modern (the default) and helvet
%      actually render on this box. Do not "upgrade" the font stack without
%      test-compiling first: kpsewhich lying about a font is the exact trap here.
%
% To get a nicer font stack + algorithm2e, install:
%   sudo apt install texlive-fonts-recommended texlive-science
% Until then, this preamble deliberately uses the default font.
%
% Usage from a document:
%   \documentclass[11pt,a4paper]{article}
%   % preamble.tex - shared house preamble for /explain-projects documents.
% Every package here is VERIFIED present in this machine's TeX Live install.
% Do not add packages without running: kpsewhich <pkg>.sty
% Notably ABSENT (do not \usepackage them, the build will fail):
%   algorithm2e, algorithm, algorithmic, algpseudocode  (texlive-science is not installed)
%   -> use the \begin{pseudocode} tcolorbox environment defined below instead.
%   lmodern, charter, mathptmx, times, mathpazo, courier, newtx*, tgtermes
%   -> their .sty files may resolve but the FONT FILES are absent, so the build
%      dies with "I can't find file `bchr8t'" or similar. Verified 2026-07-16 by
%      test-compiling each one. Only Computer Modern (the default) and helvet
%      actually render on this box. Do not "upgrade" the font stack without
%      test-compiling first: kpsewhich lying about a font is the exact trap here.
%
% To get a nicer font stack + algorithm2e, install:
%   sudo apt install texlive-fonts-recommended texlive-science
% Until then, this preamble deliberately uses the default font.
%
% Usage from a document:
%   \documentclass[11pt,a4paper]{article}
%   % preamble.tex - shared house preamble for /explain-projects documents.
% Every package here is VERIFIED present in this machine's TeX Live install.
% Do not add packages without running: kpsewhich <pkg>.sty
% Notably ABSENT (do not \usepackage them, the build will fail):
%   algorithm2e, algorithm, algorithmic, algpseudocode  (texlive-science is not installed)
%   -> use the \begin{pseudocode} tcolorbox environment defined below instead.
%   lmodern, charter, mathptmx, times, mathpazo, courier, newtx*, tgtermes
%   -> their .sty files may resolve but the FONT FILES are absent, so the build
%      dies with "I can't find file `bchr8t'" or similar. Verified 2026-07-16 by
%      test-compiling each one. Only Computer Modern (the default) and helvet
%      actually render on this box. Do not "upgrade" the font stack without
%      test-compiling first: kpsewhich lying about a font is the exact trap here.
%
% To get a nicer font stack + algorithm2e, install:
%   sudo apt install texlive-fonts-recommended texlive-science
% Until then, this preamble deliberately uses the default font.
%
% Usage from a document:
%   \documentclass[11pt,a4paper]{article}
%   \input{preamble}          % this file is copied next to the .tex

\usepackage[T1]{fontenc}
\usepackage[utf8]{inputenc}
% Font: Computer Modern (default). See the note above before changing this.
\usepackage[a4paper,margin=2.4cm,headheight=14pt]{geometry}
% microtype WITHOUT font expansion. This is load-bearing, do not "restore" it:
% cm-super is not installed, so T1 Computer Modern resolves to BITMAP fonts, and
% pdfTeX's auto-expansion needs scalable ones. With expansion on, any document
% long enough to reach a second page dies with
%   "pdfTeX error (font expansion): auto expansion is only possible with scalable fonts"
% Protrusion (the bigger visual win) still works and stays on.
% `sudo apt install cm-super` would lift this; until then, leave expansion off.
\usepackage[protrusion=true,expansion=false]{microtype}
\usepackage{graphicx}
\usepackage{xcolor}
\usepackage{booktabs}
\usepackage{enumitem}
\usepackage{fancyhdr}
\usepackage{titlesec}
\usepackage{tcolorbox}
\usepackage{listings}
\usepackage{float}
\usepackage{caption}
\usepackage{tikz}
\usepackage{forest}
\usepackage{pgfplots}
\pgfplotsset{compat=1.18}
\usetikzlibrary{arrows.meta,positioning,shapes.geometric,fit,backgrounds,calc,decorations.pathreplacing}
\tcbuselibrary{skins,breakable}
\usepackage[hidelinks,bookmarks=true,bookmarksnumbered=true]{hyperref}

% ===== palette ===========================================================
\definecolor{inkblue}{HTML}{1F3A5F}
\definecolor{accent}{HTML}{2E7DAF}
\definecolor{softgrey}{HTML}{F4F6F8}
\definecolor{linegrey}{HTML}{C8D0D8}
\definecolor{codebg}{HTML}{F7F7F4}
\definecolor{warnred}{HTML}{A33B3B}
\definecolor{okgreen}{HTML}{2E7D53}

% ===== headings ==========================================================
\titleformat{\section}{\Large\bfseries\color{inkblue}}{\thesection}{0.7em}{}
\titleformat{\subsection}{\large\bfseries\color{inkblue}}{\thesubsection}{0.6em}{}
\titleformat{\subsubsection}{\normalsize\bfseries\color{accent}}{\thesubsubsection}{0.5em}{}

\pagestyle{fancy}
\fancyhf{}
\fancyhead[L]{\small\color{inkblue}\nouppercase{\leftmark}}
\fancyhead[R]{\small\color{inkblue}\thepage}
\renewcommand{\headrulewidth}{0.4pt}
\renewcommand{\headrule}{\hbox to\headwidth{\color{linegrey}\leaders\hrule height \headrulewidth\hfill}}

% ===== code ==============================================================
\lstdefinestyle{house}{
  backgroundcolor=\color{codebg},
  basicstyle=\ttfamily\footnotesize,
  keywordstyle=\color{inkblue}\bfseries,
  commentstyle=\color{gray}\itshape,
  stringstyle=\color{okgreen},
  numbers=left, numberstyle=\tiny\color{gray}, numbersep=8pt,
  frame=single, rulecolor=\color{linegrey},
  breaklines=true, breakatwhitespace=false,
  showstringspaces=false, tabsize=2,
  captionpos=b, xleftmargin=14pt, framexleftmargin=14pt,
  columns=fullflexible, keepspaces=true,
  upquote=true  % render ' and " as straight quotes, not curly
}
\lstset{style=house}

% ===== callout boxes =====================================================
% \begin{keyidea}{Title} ... \end{keyidea}   - a concept a beginner must grasp
\newtcolorbox{keyidea}[1]{
  enhanced, breakable, colback=softgrey, colframe=accent,
  boxrule=0.4pt, arc=2pt, left=8pt, right=8pt, top=6pt, bottom=6pt,
  title={\bfseries #1}, coltitle=white, colbacktitle=accent, fonttitle=\small
}

% \begin{jargon}{Term} definition \end{jargon}  - define a term at first use
\newtcolorbox{jargon}[1]{
  enhanced, breakable, colback=white, colframe=linegrey,
  boxrule=0.4pt, arc=2pt, left=8pt, right=8pt, top=6pt, bottom=6pt,
  title={\bfseries Jargon: #1}, coltitle=inkblue, colbacktitle=softgrey, fonttitle=\small
}

% \begin{honest}{Title} ... \end{honest}  - brutally honest finding (rule 18)
\newtcolorbox{honest}[1]{
  enhanced, breakable, colback=white, colframe=warnred,
  boxrule=0.6pt, arc=2pt, left=8pt, right=8pt, top=6pt, bottom=6pt,
  title={\bfseries #1}, coltitle=white, colbacktitle=warnred, fonttitle=\small
}

% \begin{provenance}{path/to/file} ... \end{provenance}
% Where a file came from: hand-written, generated by a tool, scaffolded by a
% framework, copied or adapted from an upstream project, ported from an older
% module. State the EVIDENCE (a header comment, a git log entry, an upstream
% URL, a generator command) and label the confidence. If no evidence exists,
% say so in the box. Never guess a lineage into existence.
\newtcolorbox{provenance}[1]{
  enhanced, breakable, colback=white, colframe=okgreen,
  boxrule=0.4pt, arc=2pt, left=8pt, right=8pt, top=6pt, bottom=6pt,
  title={\bfseries Where this came from: #1}, coltitle=white,
  colbacktitle=okgreen, fonttitle=\small
}

% \begin{pseudocode}{Caption} ... \end{pseudocode}
% Stands in for algorithm2e, which is NOT installed. Use \Step / \Indent inside.
\newtcolorbox{pseudocode}[1]{
  enhanced, breakable, colback=codebg, colframe=inkblue,
  boxrule=0.4pt, arc=2pt, left=10pt, right=8pt, top=6pt, bottom=6pt,
  fontupper=\ttfamily\small,
  title={\bfseries\rmfamily #1}, coltitle=white, colbacktitle=inkblue, fonttitle=\small
}
\newcommand{\Step}[1]{\par\noindent #1}
\newcommand{\Indent}[1]{\par\noindent\hspace*{2em}#1}

% ===== tikz house styles =================================================
\tikzset{
  box/.style   = {rectangle, rounded corners=2pt, draw=inkblue, fill=softgrey,
                  text=inkblue, align=center, inner sep=6pt, minimum height=9mm,
                  font=\small},
  extbox/.style= {rectangle, rounded corners=2pt, draw=linegrey, fill=white,
                  text=black!70, align=center, inner sep=6pt, minimum height=9mm,
                  font=\small, dashed},
  store/.style = {cylinder, shape border rotate=90, aspect=0.25, draw=accent,
                  fill=white, text=inkblue, align=center, inner sep=5pt,
                  font=\small},
  dec/.style   = {diamond, aspect=1.6, draw=accent, fill=white, text=inkblue,
                  align=center, inner sep=2pt, font=\scriptsize},
  flow/.style  = {-{Stealth[length=2.4mm]}, draw=inkblue, thick},
  dflow/.style = {-{Stealth[length=2.4mm]}, draw=accent, thick, dashed},
  % align=center is load-bearing: without it, a \\ inside an edge label throws
  % the misleading "perhaps a missing \item" error.
  lbl/.style   = {font=\scriptsize\itshape, text=black!65, fill=white,
                  inner sep=2pt, align=center}
}

% forest default for directory trees
\forestset{
  dirtree/.style={
    for tree={
      font=\ttfamily\small, grow'=0, child anchor=west, parent anchor=south,
      anchor=west, calign=first, inner sep=2pt,
      edge path={\noexpand\path [draw, \forestoption{edge}]
        (!u.south west) +(9pt,0) |- (.child anchor) \forestoption{edge label};},
      before typesetting nodes={if n=1{insert before={[,phantom]}}{}},
      fit=band, before computing xy={l=15pt}
    }
  }
}

% ===== misc ==============================================================
\captionsetup{font=small, labelfont={bf,color=inkblue}}
\setlist[itemize]{leftmargin=1.2em, itemsep=2pt, topsep=3pt}
\setlist[enumerate]{leftmargin=1.4em, itemsep=2pt, topsep=3pt}
\newcommand{\file}[1]{\texttt{\small #1}}
\newcommand{\code}[1]{\texttt{\small #1}}
\setcounter{tocdepth}{2}
\setcounter{secnumdepth}{3}
          % this file is copied next to the .tex

\usepackage[T1]{fontenc}
\usepackage[utf8]{inputenc}
% Font: Computer Modern (default). See the note above before changing this.
\usepackage[a4paper,margin=2.4cm,headheight=14pt]{geometry}
% microtype WITHOUT font expansion. This is load-bearing, do not "restore" it:
% cm-super is not installed, so T1 Computer Modern resolves to BITMAP fonts, and
% pdfTeX's auto-expansion needs scalable ones. With expansion on, any document
% long enough to reach a second page dies with
%   "pdfTeX error (font expansion): auto expansion is only possible with scalable fonts"
% Protrusion (the bigger visual win) still works and stays on.
% `sudo apt install cm-super` would lift this; until then, leave expansion off.
\usepackage[protrusion=true,expansion=false]{microtype}
\usepackage{graphicx}
\usepackage{xcolor}
\usepackage{booktabs}
\usepackage{enumitem}
\usepackage{fancyhdr}
\usepackage{titlesec}
\usepackage{tcolorbox}
\usepackage{listings}
\usepackage{float}
\usepackage{caption}
\usepackage{tikz}
\usepackage{forest}
\usepackage{pgfplots}
\pgfplotsset{compat=1.18}
\usetikzlibrary{arrows.meta,positioning,shapes.geometric,fit,backgrounds,calc,decorations.pathreplacing}
\tcbuselibrary{skins,breakable}
\usepackage[hidelinks,bookmarks=true,bookmarksnumbered=true]{hyperref}

% ===== palette ===========================================================
\definecolor{inkblue}{HTML}{1F3A5F}
\definecolor{accent}{HTML}{2E7DAF}
\definecolor{softgrey}{HTML}{F4F6F8}
\definecolor{linegrey}{HTML}{C8D0D8}
\definecolor{codebg}{HTML}{F7F7F4}
\definecolor{warnred}{HTML}{A33B3B}
\definecolor{okgreen}{HTML}{2E7D53}

% ===== headings ==========================================================
\titleformat{\section}{\Large\bfseries\color{inkblue}}{\thesection}{0.7em}{}
\titleformat{\subsection}{\large\bfseries\color{inkblue}}{\thesubsection}{0.6em}{}
\titleformat{\subsubsection}{\normalsize\bfseries\color{accent}}{\thesubsubsection}{0.5em}{}

\pagestyle{fancy}
\fancyhf{}
\fancyhead[L]{\small\color{inkblue}\nouppercase{\leftmark}}
\fancyhead[R]{\small\color{inkblue}\thepage}
\renewcommand{\headrulewidth}{0.4pt}
\renewcommand{\headrule}{\hbox to\headwidth{\color{linegrey}\leaders\hrule height \headrulewidth\hfill}}

% ===== code ==============================================================
\lstdefinestyle{house}{
  backgroundcolor=\color{codebg},
  basicstyle=\ttfamily\footnotesize,
  keywordstyle=\color{inkblue}\bfseries,
  commentstyle=\color{gray}\itshape,
  stringstyle=\color{okgreen},
  numbers=left, numberstyle=\tiny\color{gray}, numbersep=8pt,
  frame=single, rulecolor=\color{linegrey},
  breaklines=true, breakatwhitespace=false,
  showstringspaces=false, tabsize=2,
  captionpos=b, xleftmargin=14pt, framexleftmargin=14pt,
  columns=fullflexible, keepspaces=true,
  upquote=true  % render ' and " as straight quotes, not curly
}
\lstset{style=house}

% ===== callout boxes =====================================================
% \begin{keyidea}{Title} ... \end{keyidea}   - a concept a beginner must grasp
\newtcolorbox{keyidea}[1]{
  enhanced, breakable, colback=softgrey, colframe=accent,
  boxrule=0.4pt, arc=2pt, left=8pt, right=8pt, top=6pt, bottom=6pt,
  title={\bfseries #1}, coltitle=white, colbacktitle=accent, fonttitle=\small
}

% \begin{jargon}{Term} definition \end{jargon}  - define a term at first use
\newtcolorbox{jargon}[1]{
  enhanced, breakable, colback=white, colframe=linegrey,
  boxrule=0.4pt, arc=2pt, left=8pt, right=8pt, top=6pt, bottom=6pt,
  title={\bfseries Jargon: #1}, coltitle=inkblue, colbacktitle=softgrey, fonttitle=\small
}

% \begin{honest}{Title} ... \end{honest}  - brutally honest finding (rule 18)
\newtcolorbox{honest}[1]{
  enhanced, breakable, colback=white, colframe=warnred,
  boxrule=0.6pt, arc=2pt, left=8pt, right=8pt, top=6pt, bottom=6pt,
  title={\bfseries #1}, coltitle=white, colbacktitle=warnred, fonttitle=\small
}

% \begin{provenance}{path/to/file} ... \end{provenance}
% Where a file came from: hand-written, generated by a tool, scaffolded by a
% framework, copied or adapted from an upstream project, ported from an older
% module. State the EVIDENCE (a header comment, a git log entry, an upstream
% URL, a generator command) and label the confidence. If no evidence exists,
% say so in the box. Never guess a lineage into existence.
\newtcolorbox{provenance}[1]{
  enhanced, breakable, colback=white, colframe=okgreen,
  boxrule=0.4pt, arc=2pt, left=8pt, right=8pt, top=6pt, bottom=6pt,
  title={\bfseries Where this came from: #1}, coltitle=white,
  colbacktitle=okgreen, fonttitle=\small
}

% \begin{pseudocode}{Caption} ... \end{pseudocode}
% Stands in for algorithm2e, which is NOT installed. Use \Step / \Indent inside.
\newtcolorbox{pseudocode}[1]{
  enhanced, breakable, colback=codebg, colframe=inkblue,
  boxrule=0.4pt, arc=2pt, left=10pt, right=8pt, top=6pt, bottom=6pt,
  fontupper=\ttfamily\small,
  title={\bfseries\rmfamily #1}, coltitle=white, colbacktitle=inkblue, fonttitle=\small
}
\newcommand{\Step}[1]{\par\noindent #1}
\newcommand{\Indent}[1]{\par\noindent\hspace*{2em}#1}

% ===== tikz house styles =================================================
\tikzset{
  box/.style   = {rectangle, rounded corners=2pt, draw=inkblue, fill=softgrey,
                  text=inkblue, align=center, inner sep=6pt, minimum height=9mm,
                  font=\small},
  extbox/.style= {rectangle, rounded corners=2pt, draw=linegrey, fill=white,
                  text=black!70, align=center, inner sep=6pt, minimum height=9mm,
                  font=\small, dashed},
  store/.style = {cylinder, shape border rotate=90, aspect=0.25, draw=accent,
                  fill=white, text=inkblue, align=center, inner sep=5pt,
                  font=\small},
  dec/.style   = {diamond, aspect=1.6, draw=accent, fill=white, text=inkblue,
                  align=center, inner sep=2pt, font=\scriptsize},
  flow/.style  = {-{Stealth[length=2.4mm]}, draw=inkblue, thick},
  dflow/.style = {-{Stealth[length=2.4mm]}, draw=accent, thick, dashed},
  % align=center is load-bearing: without it, a \\ inside an edge label throws
  % the misleading "perhaps a missing \item" error.
  lbl/.style   = {font=\scriptsize\itshape, text=black!65, fill=white,
                  inner sep=2pt, align=center}
}

% forest default for directory trees
\forestset{
  dirtree/.style={
    for tree={
      font=\ttfamily\small, grow'=0, child anchor=west, parent anchor=south,
      anchor=west, calign=first, inner sep=2pt,
      edge path={\noexpand\path [draw, \forestoption{edge}]
        (!u.south west) +(9pt,0) |- (.child anchor) \forestoption{edge label};},
      before typesetting nodes={if n=1{insert before={[,phantom]}}{}},
      fit=band, before computing xy={l=15pt}
    }
  }
}

% ===== misc ==============================================================
\captionsetup{font=small, labelfont={bf,color=inkblue}}
\setlist[itemize]{leftmargin=1.2em, itemsep=2pt, topsep=3pt}
\setlist[enumerate]{leftmargin=1.4em, itemsep=2pt, topsep=3pt}
\newcommand{\file}[1]{\texttt{\small #1}}
\newcommand{\code}[1]{\texttt{\small #1}}
\setcounter{tocdepth}{2}
\setcounter{secnumdepth}{3}
          % this file is copied next to the .tex

\usepackage[T1]{fontenc}
\usepackage[utf8]{inputenc}
% Font: Computer Modern (default). See the note above before changing this.
\usepackage[a4paper,margin=2.4cm,headheight=14pt]{geometry}
% microtype WITHOUT font expansion. This is load-bearing, do not "restore" it:
% cm-super is not installed, so T1 Computer Modern resolves to BITMAP fonts, and
% pdfTeX's auto-expansion needs scalable ones. With expansion on, any document
% long enough to reach a second page dies with
%   "pdfTeX error (font expansion): auto expansion is only possible with scalable fonts"
% Protrusion (the bigger visual win) still works and stays on.
% `sudo apt install cm-super` would lift this; until then, leave expansion off.
\usepackage[protrusion=true,expansion=false]{microtype}
\usepackage{graphicx}
\usepackage{xcolor}
\usepackage{booktabs}
\usepackage{enumitem}
\usepackage{fancyhdr}
\usepackage{titlesec}
\usepackage{tcolorbox}
\usepackage{listings}
\usepackage{float}
\usepackage{caption}
\usepackage{tikz}
\usepackage{forest}
\usepackage{pgfplots}
\pgfplotsset{compat=1.18}
\usetikzlibrary{arrows.meta,positioning,shapes.geometric,fit,backgrounds,calc,decorations.pathreplacing}
\tcbuselibrary{skins,breakable}
\usepackage[hidelinks,bookmarks=true,bookmarksnumbered=true]{hyperref}

% ===== palette ===========================================================
\definecolor{inkblue}{HTML}{1F3A5F}
\definecolor{accent}{HTML}{2E7DAF}
\definecolor{softgrey}{HTML}{F4F6F8}
\definecolor{linegrey}{HTML}{C8D0D8}
\definecolor{codebg}{HTML}{F7F7F4}
\definecolor{warnred}{HTML}{A33B3B}
\definecolor{okgreen}{HTML}{2E7D53}

% ===== headings ==========================================================
\titleformat{\section}{\Large\bfseries\color{inkblue}}{\thesection}{0.7em}{}
\titleformat{\subsection}{\large\bfseries\color{inkblue}}{\thesubsection}{0.6em}{}
\titleformat{\subsubsection}{\normalsize\bfseries\color{accent}}{\thesubsubsection}{0.5em}{}

\pagestyle{fancy}
\fancyhf{}
\fancyhead[L]{\small\color{inkblue}\nouppercase{\leftmark}}
\fancyhead[R]{\small\color{inkblue}\thepage}
\renewcommand{\headrulewidth}{0.4pt}
\renewcommand{\headrule}{\hbox to\headwidth{\color{linegrey}\leaders\hrule height \headrulewidth\hfill}}

% ===== code ==============================================================
\lstdefinestyle{house}{
  backgroundcolor=\color{codebg},
  basicstyle=\ttfamily\footnotesize,
  keywordstyle=\color{inkblue}\bfseries,
  commentstyle=\color{gray}\itshape,
  stringstyle=\color{okgreen},
  numbers=left, numberstyle=\tiny\color{gray}, numbersep=8pt,
  frame=single, rulecolor=\color{linegrey},
  breaklines=true, breakatwhitespace=false,
  showstringspaces=false, tabsize=2,
  captionpos=b, xleftmargin=14pt, framexleftmargin=14pt,
  columns=fullflexible, keepspaces=true,
  upquote=true  % render ' and " as straight quotes, not curly
}
\lstset{style=house}

% ===== callout boxes =====================================================
% \begin{keyidea}{Title} ... \end{keyidea}   - a concept a beginner must grasp
\newtcolorbox{keyidea}[1]{
  enhanced, breakable, colback=softgrey, colframe=accent,
  boxrule=0.4pt, arc=2pt, left=8pt, right=8pt, top=6pt, bottom=6pt,
  title={\bfseries #1}, coltitle=white, colbacktitle=accent, fonttitle=\small
}

% \begin{jargon}{Term} definition \end{jargon}  - define a term at first use
\newtcolorbox{jargon}[1]{
  enhanced, breakable, colback=white, colframe=linegrey,
  boxrule=0.4pt, arc=2pt, left=8pt, right=8pt, top=6pt, bottom=6pt,
  title={\bfseries Jargon: #1}, coltitle=inkblue, colbacktitle=softgrey, fonttitle=\small
}

% \begin{honest}{Title} ... \end{honest}  - brutally honest finding (rule 18)
\newtcolorbox{honest}[1]{
  enhanced, breakable, colback=white, colframe=warnred,
  boxrule=0.6pt, arc=2pt, left=8pt, right=8pt, top=6pt, bottom=6pt,
  title={\bfseries #1}, coltitle=white, colbacktitle=warnred, fonttitle=\small
}

% \begin{provenance}{path/to/file} ... \end{provenance}
% Where a file came from: hand-written, generated by a tool, scaffolded by a
% framework, copied or adapted from an upstream project, ported from an older
% module. State the EVIDENCE (a header comment, a git log entry, an upstream
% URL, a generator command) and label the confidence. If no evidence exists,
% say so in the box. Never guess a lineage into existence.
\newtcolorbox{provenance}[1]{
  enhanced, breakable, colback=white, colframe=okgreen,
  boxrule=0.4pt, arc=2pt, left=8pt, right=8pt, top=6pt, bottom=6pt,
  title={\bfseries Where this came from: #1}, coltitle=white,
  colbacktitle=okgreen, fonttitle=\small
}

% \begin{pseudocode}{Caption} ... \end{pseudocode}
% Stands in for algorithm2e, which is NOT installed. Use \Step / \Indent inside.
\newtcolorbox{pseudocode}[1]{
  enhanced, breakable, colback=codebg, colframe=inkblue,
  boxrule=0.4pt, arc=2pt, left=10pt, right=8pt, top=6pt, bottom=6pt,
  fontupper=\ttfamily\small,
  title={\bfseries\rmfamily #1}, coltitle=white, colbacktitle=inkblue, fonttitle=\small
}
\newcommand{\Step}[1]{\par\noindent #1}
\newcommand{\Indent}[1]{\par\noindent\hspace*{2em}#1}

% ===== tikz house styles =================================================
\tikzset{
  box/.style   = {rectangle, rounded corners=2pt, draw=inkblue, fill=softgrey,
                  text=inkblue, align=center, inner sep=6pt, minimum height=9mm,
                  font=\small},
  extbox/.style= {rectangle, rounded corners=2pt, draw=linegrey, fill=white,
                  text=black!70, align=center, inner sep=6pt, minimum height=9mm,
                  font=\small, dashed},
  store/.style = {cylinder, shape border rotate=90, aspect=0.25, draw=accent,
                  fill=white, text=inkblue, align=center, inner sep=5pt,
                  font=\small},
  dec/.style   = {diamond, aspect=1.6, draw=accent, fill=white, text=inkblue,
                  align=center, inner sep=2pt, font=\scriptsize},
  flow/.style  = {-{Stealth[length=2.4mm]}, draw=inkblue, thick},
  dflow/.style = {-{Stealth[length=2.4mm]}, draw=accent, thick, dashed},
  % align=center is load-bearing: without it, a \\ inside an edge label throws
  % the misleading "perhaps a missing \item" error.
  lbl/.style   = {font=\scriptsize\itshape, text=black!65, fill=white,
                  inner sep=2pt, align=center}
}

% forest default for directory trees
\forestset{
  dirtree/.style={
    for tree={
      font=\ttfamily\small, grow'=0, child anchor=west, parent anchor=south,
      anchor=west, calign=first, inner sep=2pt,
      edge path={\noexpand\path [draw, \forestoption{edge}]
        (!u.south west) +(9pt,0) |- (.child anchor) \forestoption{edge label};},
      before typesetting nodes={if n=1{insert before={[,phantom]}}{}},
      fit=band, before computing xy={l=15pt}
    }
  }
}

% ===== misc ==============================================================
\captionsetup{font=small, labelfont={bf,color=inkblue}}
\setlist[itemize]{leftmargin=1.2em, itemsep=2pt, topsep=3pt}
\setlist[enumerate]{leftmargin=1.4em, itemsep=2pt, topsep=3pt}
\newcommand{\file}[1]{\texttt{\small #1}}
\newcommand{\code}[1]{\texttt{\small #1}}
\setcounter{tocdepth}{2}
\setcounter{secnumdepth}{3}


\DeclareUnicodeCharacter{2500}{-}
\DeclareUnicodeCharacter{2550}{=}
\DeclareUnicodeCharacter{2764}{\ensuremath{\heartsuit}}
\lstset{extendedchars=true,
  literate={·}{{\textperiodcentered}}1
           {─}{{-}}1
           {═}{{=}}1
           {❤}{{\ensuremath{\heartsuit}}}1
}

\title{AI Digitalise\\[2mm]\large Brief: the whole project in one sitting}
\author{Digitalise Agency}
\date{2026-08-11}

\begin{document}
\maketitle
\tableofcontents
\newpage

% ===========================================================================
\section{What it is, and the one thing to know first}
% ===========================================================================

This repository holds a marketing web page for \textbf{ai.digitalise.agency}: one
page, dark navy, whose job is to convince a visitor that Digitalise Agency can
build AI systems for them and then get them to book a call.

Six files, 1{,}591 lines, no framework, no build step, no package manager, no
database and no server code. That is unusual enough in 2026 to be worth stating
as a decision rather than an absence.

\begin{keyidea}{This is not the page being served at that address}
The design here was finished, then deliberately parked. On 2026-07-29 a single
commit moved every file into a folder called \file{new-design-reference/} and
renamed the publishing workflow so it could no longer run. Its message says why:
the live site is a different build, produced elsewhere.

Fetching \code{https://ai.digitalise.agency/} on 2026-08-11 confirms it. The
address answers with \code{200} and serves a page with a different title
entirely. Read this repository as a design reference, never as the live site.
\end{keyidea}

\begin{figure}[H]\centering
\begin{tikzpicture}[node distance=10mm and 18mm]
  \node[box]                        (v)  {A visitor\\\scriptsize from a link or an ad};
  \node[box, right=of v]            (p)  {The page\\\scriptsize \file{index.html}};
  \node[extbox, right=of p]         (cal){cal.com\\\scriptsize booking calendar};
  \node[box, below=of p]            (out){A booked call\\\scriptsize the only outcome that counts};
  \draw[flow]  (v) -- node[lbl]{opens} (p);
  \draw[dflow] (p) -- node[lbl]{Book a call} (cal);
  \draw[flow]  (cal) |- node[lbl, pos=0.25]{picks a time} (out);
\end{tikzpicture}
\caption{The entire purpose. Dashed boxes are services the agency does not run.}
\end{figure}

\begin{jargon}{Repository}
A folder whose complete change history is recorded by a program called git, so
every edit can be listed, inspected and undone.
\end{jargon}

% ===========================================================================
\section{The mental model}
% ===========================================================================

A web page is made of three kinds of file. This project has exactly one of each,
plus three files that are about hosting rather than about the page.

\begin{figure}[H]\centering
\begin{tikzpicture}[node distance=11mm and 14mm]
  \node[box]                 (html) {\file{index.html}\\\scriptsize 597 lines: the words and structure};
  \node[box, right=of html]  (css)  {\file{css/site.css}\\\scriptsize 768 lines: what it looks like};
  \node[box, right=of css]   (js)   {\file{js/main.js}\\\scriptsize 79 lines: what moves};
  \node[extbox, below=of css](br)   {The visitor's browser\\\scriptsize assembles all three};
  \draw[flow] (html) -- node[lbl]{names} (css);
  \draw[flow] (css)  -- node[lbl]{then} (js);
  \draw[flow] (html) |- (br);
  \draw[flow] (css)  -- (br);
  \draw[flow] (js)   |- (br);
\end{tikzpicture}
\caption{Nothing runs on a server. Everything happens on the visitor's machine.}
\end{figure}

\begin{keyidea}{The arrows only point one way}
\file{index.html} names the other two files. Neither knows anything about it: the
stylesheet styles class names it hopes exist, and the JavaScript looks for
elements that may not be there, with a guard at every lookup so it does nothing
quietly when they are missing.

That is why the page never crashes, and why a mistyped class name produces
something that looks wrong rather than something that reports an error. Nothing
anywhere checks that the three files agree with each other.
\end{keyidea}

% ===========================================================================
\section{How it works, end to end}
% ===========================================================================

\begin{figure}[H]\centering
\begin{tikzpicture}[node distance=8mm and 12mm]
  \node[box]             (a) {Browser downloads\\\file{index.html}};
  \node[box, right=of a] (b) {Requests fonts\\and the stylesheet};
  \node[box, right=of b] (c) {Draws the page\\top to bottom};
  \node[box, below=of c] (d) {Reaches \code{<script>}\\at the very end};
  \node[box, left=of d]  (e) {\file{main.js} runs,\\attaches watchers};
  \node[box, left=of e]  (f) {Visitor scrolls,\\watchers fire};
  \draw[flow] (a) -- (b); \draw[flow] (b) -- (c); \draw[flow] (c) -- (d);
  \draw[flow] (d) -- (e); \draw[flow] (e) -- (f);
\end{tikzpicture}
\caption{One page view, start to finish.}
\end{figure}

The page stacks nine bands: a fixed navigation bar, a full-screen hero with eight
small cards drifting behind it, four large statistics, a five-step pipeline
diagram, six portfolio cards, three client result cards, six service tiles, a
booking panel, and a one-line footer.

Three things then react to scrolling, and that is the whole of the interactivity:

\begin{enumerate}
\item the navigation bar grows a frosted dark background past 40 pixels of
  scroll;
\item the four statistics count up from zero over 2.2 seconds when they come into
  view;
\item most blocks fade upward into place as they arrive.
\end{enumerate}

There is no form, nothing is saved, and nothing is sent anywhere. The only
outbound action a visitor can take is following the booking link, which opens a
cal.com calendar in a new tab.

\begin{honest}{Nothing measures whether the page works}
No analytics of any kind: no counter, no tag manager, no event when somebody
presses ``Book a call''. The one question the page exists to answer cannot be
answered from the page. Whether that is a privacy choice or an oversight is not
recorded anywhere, so this document does not guess.
\end{honest}

% ===========================================================================
\section{The main components}
% ===========================================================================

\begin{table}[H]\centering\small
\begin{tabular}{@{}llrp{5.6cm}@{}}
\toprule
File & Kind & Lines & What it does \\
\midrule
\file{new-design-reference/index.html} & HTML & 597 & Every word on the page, and its structure. Holds all the icons as inline drawings \\
\file{new-design-reference/css/site.css} & CSS & 768 & All appearance, layout, movement, and the two mobile layouts \\
\file{new-design-reference/js/main.js} & JavaScript & 79 & The three scroll behaviours, with no dependencies at all \\
\file{new-design-reference/netlify.toml} & config & 11 & Hosting settings for a host this project does not use \\
\file{new-design-reference/assets/ASSETS.md} & notes & 33 & A checklist of twelve assets, none of which were ever produced \\
\file{.github/workflows-disabled/deploy.yml.disabled} & workflow & 103 & The publishing pipeline, switched off \\
\bottomrule
\end{tabular}
\caption{Every file git tracks. The Deep Dive walks all six line by line, in
sections 7.1 to 7.6.}
\end{table}

\subsection{\file{index.html}: the page}

The content and the structure. It is unusually self-contained: every icon is an
SVG drawing written directly into the file rather than loaded as an image, so the
page needs no image downloads at all.

Two of its habits are worth knowing before you edit it. First, appearance is
sometimes written directly onto elements with \code{style} attributes rather than
left to the stylesheet, and an inline style always wins, which is how one
stylesheet rule ends up permanently overridden. Second, the four statistics keep
their real values in \code{data-} attributes and display the character \code{0}
until JavaScript replaces it. Deep Dive section 7.1.

\subsection{\file{site.css}: the design system}

The most valuable file in the repository, and the one worth borrowing. It opens
with twenty named values (the ``tokens'') covering surfaces, text, accents,
borders, typefaces, radii, a motion curve and a content width. Its own second line
records where the scheme came from: the tokens mirror those of the Gantry project
elsewhere in this workspace.

\begin{keyidea}{Why the token names are numbers rather than colours}
Calling a colour \code{--s1} instead of \code{--dark-navy} names what it is
\emph{for} (the surface one step up from the page) rather than how it looks.
Redesign in a light theme and \code{--s1} still means the right thing while
\code{--dark-navy} becomes a lie. This is the single most transferable idea in the
repository.
\end{keyidea}

It ends with the best thing in the project: a block that honours the visitor's
reduced-motion preference, crushing every animation to nothing \emph{and} forcing
the fade-in elements to full opacity so they still appear. That second half is the
detail most implementations get wrong. Deep Dive section 7.2.

\subsection{\file{main.js}: the three behaviours}

Seventy-nine lines, no dependencies, and better written than its size suggests: a
single wrapper keeping its names private, a guard at every lookup, a small shared
helper built on \code{IntersectionObserver} that four one-line calls then reuse,
and correct use of the flags that keep scrolling smooth and stop each animation
firing twice.

Its interface with the rest of the project is seven strings: three selectors, two
class names it adds, and two more it looks up. Break any one of them and the page
keeps working, minus one behaviour, silently. Deep Dive section 7.3.

\subsection{The three files about hosting}

\file{netlify.toml} declares four sensible security headers for a hosting service
this project does not use. \file{ASSETS.md} lists the twelve assets needed to
finish the design, none of which exist. The workflow publishes the site to the
agency's server when a commit subject begins with \code{[deploy]}, and it is
switched off. Deep Dive sections 7.4 to 7.6.

% ===========================================================================
\section{How it is operated}
% ===========================================================================

\subsection{Looking at it locally}

\begin{lstlisting}[language=bash,caption={Serve it, do not open the file directly}]
cd new-design-reference
python3 -m http.server 8000
# then open http://localhost:8000/
\end{lstlisting}

Opening \file{index.html} straight off the disk renders it completely unstyled,
which looks like catastrophic breakage and is not. The page asks for its
stylesheet at \code{/css/site.css}, and the leading slash means ``from the top of
the site'', which off a disk is the root of your hard drive. This one character is
the most likely thing to make a newcomer think the project is broken.

\subsection{Publishing}

\begin{figure}[H]\centering
\begin{tikzpicture}[node distance=9mm and 13mm]
  \node[box]                (push) {Push to \code{main}};
  \node[dec, right=of push] (chk)  {subject starts\\\code{[deploy]}?};
  \node[box, right=of chk]  (ssh)  {Install key,\\fetch host key};
  \node[box, below=of ssh]  (sync) {\code{rsync} to the VPS};
  \node[box, left=of sync]  (ver)  {Check the site\\still answers};
  \node[box, left=of ver]   (stop) {Do nothing};
  \draw[flow] (push) -- (chk);
  \draw[flow] (chk) -- node[lbl]{yes} (ssh);
  \draw[flow] (chk) |- node[lbl,pos=0.3]{no} (stop);
  \draw[flow] (ssh) -- (sync);
  \draw[flow] (sync) -- (ver);
\end{tikzpicture}
\caption{The publishing pipeline as written. Drawn for reference: it is switched
off, twice over, and has been since 2026-07-29.}
\end{figure}

The \code{[deploy]} marker must \emph{begin} the commit subject line. Writing
``fix the header and [deploy]'' does nothing at all. This convention is shared by
every repository in this workspace.

The workflow is disabled two ways: its folder is named
\file{.github/workflows-disabled} rather than \file{.github/workflows}, and the
file ends in \file{.disabled} rather than \file{.yml}. Either alone would be
enough, which tells you the intent was not to pause it briefly.

\subsection{Claude Code tooling}

\textbf{This project ships none}, which was checked rather than assumed: there is
no \file{.claude/} folder, no \file{CLAUDE.md}, no skill and no hook.

There is no README either, and that matters more than it sounds. The central fact
about this repository, that it is parked and the live site is elsewhere, exists
only in a commit message, invisible unless somebody thinks to run \code{git log}.

The two PDFs in \file{docs/explainer/} are produced by a workspace-level command,
\code{/explain-projects}, which lives outside this repository.

% ===========================================================================
\section{Configuration and deployment at a glance}
% ===========================================================================

Nothing is configurable at run time: no settings file, no environment file. Three
values behave like configuration and are written into the source:

\begin{table}[H]\centering\small
\begin{tabular}{@{}lll@{}}
\toprule
Value & Where & Controls \\
\midrule
\code{https://cal.com/davonex/30min} & \file{index.html} line 565 & which calendar opens \\
\code{hello@digitalise.agency} & \file{index.html} line 590 & the footer contact \\
The twenty colour tokens & \file{site.css} lines 5--25 & the entire colour scheme \\
\bottomrule
\end{tabular}
\caption{Values that behave like settings.}
\end{table}

The repository holds no passwords, keys or tokens, and it should not: a page that
runs entirely in a visitor's browser can keep no secret. Three secrets are stored
in GitHub and referenced by the workflow (\code{VPS\_HOST}, \code{VPS\_USER},
\code{VPS\_SSH\_KEY}); the server's IP address is also written into a comment in
that same file, which defeats the point of storing it as a secret.

The page has exactly two run-time dependencies, both reached by the visitor's
browser: Google Fonts for two typefaces, and cal.com for the booking calendar.

% ===========================================================================
\section{Honest findings}
% ===========================================================================

The Deep Dive argues each of these at the line where it arises. Summarized:

\subsection{Correctness}

\begin{itemize}
\item \textbf{The live site's canonical link points at another domain.} The page
  currently served at ai.digitalise.agency declares \code{salmanadnan.com} as its
  canonical URL, which asks search engines to index that site instead. Verified
  live on 2026-08-11. It cannot be fixed from this repository, and whoever
  maintains the live build should know.
\item \textbf{Restoring the deploy workflow unchanged would break the live site.}
  It copies the repository root, which since the move contains a folder rather
  than a home page, and it copies with \code{--delete}.
\item \textbf{A failed post-deploy check is only a warning}, so the workflow would
  finish green even if the site never came back after an \code{rsync --delete}.
\end{itemize}

\subsection{Unfinished work}

\begin{itemize}
\item \textbf{None of the twelve assets in \file{ASSETS.md} exist}, including the
  favicon and the social preview image. Portfolio thumbnails are still colour
  gradients and client logos are still lettered circles.
\item \textbf{Five ``UPDATE'' comments} remain in \file{index.html} marking work
  that may or may not have been done, and two of the three instructions in
  \file{ASSETS.md} refer to code that no longer exists.
\item \textbf{The booking link is a placeholder}, pointing at another project's
  calendar. The same placeholder is a known outstanding item on the agency's main
  website, so it is a shared loose end rather than a local slip.
\item \textbf{One client avatar shows the wrong initial}, a survivor of the
  placeholder set.
\end{itemize}

\subsection{Presentation}

\begin{honest}{The hero backdrop displays invented engagement numbers}
The eight drifting cards show 847K views, 234K impressions, 1.2M, 512K, 1.8M and
19K reactions. None is a measurement; they arrived fully formed in the first
commit. They are hidden from screen readers and are plainly decorative, but a
sighted visitor reads them as claims.

This is worth naming because the same author set a higher standard elsewhere in
the same repository: \file{ASSETS.md} asks that the four large statistics be
``verified from Gantry / n8n data'', meaning traceable to the agency's own
systems.
\end{honest}

Also: the six portfolio cards show a play button on hover and play nothing, and
the footer link reads \code{digitalise.agency} while pointing at this site.

\subsection{Accessibility}

\begin{itemize}
\item \textbf{The counting numbers ignore the reduced-motion preference.} The
  stylesheet honours it meticulously; the JavaScript never checks. The intent is
  demonstrably there and simply was not carried into the other file. Four lines
  fix it, given in Deep Dive section 7.3.
\item \textbf{Neither button has a focus state}, so keyboard navigation relies on
  whatever outline the browser draws, which on a dark background is often nearly
  invisible. These are the page's only interactive elements.
\item \textbf{The small blue uppercase labels are worth a contrast measurement.}
  This document does not have one and so raises it rather than asserting a
  failure.
\end{itemize}

\subsection{Dead and contradictory code}

Two colour tokens defined and never read; one stylesheet rule targeting a class no
element carries; one rule permanently overridden by inline styles; three lines of
responsive code that show a tile and then hide it again; a doubled flex layout in
the navigation; smooth scrolling implemented twice, once in CSS and once in
JavaScript, where the JavaScript copy also carries a latent crash; a class added
by JavaScript that nothing styles; and a whole hosting configuration file for a
host that is not used.

% ===========================================================================
\section{The story, in brief}
% ===========================================================================

\subsection{What it delivered}

A complete, self-contained one-page marketing site with no dependencies to
install; a hosting configuration setting four security headers; a GitHub Actions
workflow that publishes over SSH with retries and a post-deploy check; and an
asset checklist. What it did not deliver is a live site.

\subsection{When it was built}

Nine commits across two calendar days, 2026-07-28 and 2026-07-29, and nothing
since. The shape of that history is the interesting part: the first day is the
design and the second day is almost entirely deployment. Four of the five commits
on the second day fix the publishing pipeline, and the fifth switches it off.

This is calendar time, not effort. Git records when commits happened and nothing
about hours spent, and no document in the repository records effort, so none is
claimed.

\subsection{The four challenges that mattered}

Each is claimed because the repository carries a scar for it.

\begin{enumerate}
\item \textbf{SSH access to the server}, fixed across three consecutive commits:
  secrets missing, then the wrong key, then a \code{deploy} user without
  permission to traverse into \code{/root}. That last one is genuinely unobvious
  and its workaround survives in three \code{rsync} flags.
\item \textbf{GitHub's runners being refused by the server.} Two independent retry
  loops were written into a fifty-line script, and the sibling repository
  \code{salmanadnan-com} carries a commit titled ``Retry, runner IP lottery'',
  which confirms the diagnosis.
\item \textbf{Em-dashes.} One commit exists solely to remove them, which is the
  workspace house rule showing up as real work.
\item \textbf{Being overtaken.} The final commit is a retreat, not a fix: the page
  was finished and working when something else became the live site. That is the
  largest challenge here and it was not technical.
\end{enumerate}

\subsection{Where it stands}

Built, complete, responsive, and parked. Superseded by a different build at the
same address. Never finished as a product, since none of its own launch checklist
was delivered. Dormant since 2026-07-29.

\subsection{Do and do not}

\emph{From here on this is advice and judgement, not a record of decisions.}

\textbf{Do:}
\begin{itemize}
\item Read \code{git log} before touching anything: it holds the only statement of
  this repository's status.
\item Serve the folder over HTTP to view it, never open the file directly.
\item Keep the \code{[deploy]} convention and both retry loops if the workflow is
  ever restored. They were paid for.
\item Treat the stylesheet as the reusable asset. It is a coherent dark theme in
  768 readable lines with no tooling.
\end{itemize}

\textbf{Do not:}
\begin{itemize}
\item Re-enable the deploy workflow without changing the \code{rsync} source path.
  It would copy a folder instead of a page and take the live site down with it.
\item Copy the hero tiles into anything real without replacing their numbers.
\item Trust \file{ASSETS.md}. Two of its three instructions describe code that no
  longer exists.
\item Assume this repository is what visitors see. It is not, and that is the
  easiest mistake to make here.
\item Delete the repository because it is not live. The stylesheet is reusable and
  the workflow encodes three separate lessons about SSH on GitHub runners.
\end{itemize}

\subsection{What would change, starting over}

\emph{Opinion, argued from the code.} Keep all of it: the no-build-step decision,
the token block, the three-behaviour JavaScript and the reduced-motion handling
are right for a page this size, and a framework would have added a hundred times
the machinery for nothing a visitor could perceive.

Change five things. Write a README on day one saying what the site is and where it
is served from, because its absence is what makes this repository confusing. Use a
relative stylesheet path so it can be opened off a disk. Self-host the two fonts,
matching the agency's own main site. Produce the favicon and social image before
the first deploy rather than listing them. And delete a placeholder comment in the
same commit that removes the placeholder.

Drop \file{netlify.toml} entirely, or move its four security headers into the
Caddy configuration on the server so they actually take effect.

\subsection{Open threads}

Twelve assets never produced; a placeholder booking calendar shared with the
agency's main site; five stale \code{UPDATE} comments; an undocumented sibling
relationship with the \code{salmanadnan-com} repository, which shares an identical
file list and was built the same week; and the canonical link on the live site
pointing at the wrong domain.

% ===========================================================================
\section{Glossary}
% ===========================================================================

\begin{description}[leftmargin=0pt,style=nextline]
\item[ARIA] Attributes beginning \code{aria-} that describe a page to screen
  readers. They change nothing visually.
\item[Breakpoint] A screen width at which the layout changes. This project has
  two, at 900 and 600 pixels.
\item[Class] A label attached to elements so CSS and JavaScript can find them.
\item[clamp] A CSS function taking a minimum, a preferred value and a maximum, used
  here so headings scale with the window between fixed limits.
\item[CSS] Cascading Style Sheets: the file controlling appearance.
\item[Custom property] A named value in CSS, written \code{--name} and read back
  with \code{var(--name)}. The ``tokens'' are these.
\item[Data attribute] An attribute beginning \code{data-} that HTML ignores and
  JavaScript can read. The four statistics keep their real values in these.
\item[Flexbox] A layout system arranging children in a row or column.
\item[Grid] A layout system arranging children in rows and columns at once.
\item[HTML] HyperText Markup Language: the file holding content and structure.
\item[IntersectionObserver] A browser facility reporting when an element scrolls
  into view. All three scroll behaviours are built on it.
\item[JavaScript] The language running in the browser after the page loads.
\item[Media query] CSS applying only when a condition about the screen is true.
\item[prefers-reduced-motion] A system setting asking pages to avoid animation,
  usually because it causes nausea or migraine. The stylesheet honours it; the
  JavaScript does not.
\item[Pseudo-element] An extra box CSS creates with no tag in the HTML, written
  \code{::before} or \code{::after}, used here for the pipeline's flowing lines
  and the glow above the booking panel.
\item[rsync] A program that copies files between computers, transferring only what
  differs.
\item[Selector] The part of a CSS rule saying which elements it applies to.
\item[SSH] Secure Shell: the standard way to connect to another computer securely.
\item[SVG] A picture described as shapes and coordinates rather than pixels, so it
  stays sharp at any size and can sit inside the HTML file.
\item[Token] In this project, one of the twenty named values at the top of the
  stylesheet.
\item[VPS] Virtual Private Server: a rented computer the agency controls remotely.
\item[Workflow] A file describing a job GitHub runs automatically when something
  happens in the repository.
\end{description}

\vspace{1em}
\noindent
\textbf{For the full treatment}, including every line of every file with its
provenance, its connections and its failure modes, see
\file{docs/explainer/deep-dive.pdf} in this same folder.

\end{document}
